\documentclass[12pt, a4paper]{article}

\usepackage[T1]{fontenc}
\usepackage[utf8]{inputenc}
\usepackage{lmodern}         

\usepackage[margin=2.5cm]{geometry}

\usepackage{amsmath, amssymb}  
\usepackage{graphicx}          
\usepackage{booktabs}         
\usepackage{hyperref}        
%\usepackage{cite}             
\usepackage{caption}          
\usepackage{subcaption}        
\usepackage{siunitx}           
\usepackage{lineno}            
\usepackage{setspace}        
\usepackage[backend=biber, style=authoryear]{biblatex}
\addbibresource{references.bib}    

\hypersetup{
    colorlinks = true,
    linkcolor  = blue,
    citecolor  = blue,
    urlcolor   = blue
}



\title{The Unified Hierarchical Reasoning Controller}

\author{
    Christon Nunes\textsuperscript{1} \and
    Cade Coker\textsuperscript{2} \and
    Dr. Craig Ramlal\textsuperscript{3}
}

\date{
    \textsuperscript{1}Department of Electrical and Computer Engineering, University of the West Indies St.Augustine \\
    \small Corresponding Author: \href{mailto:nuneschriston20@gmail.com}{nuneschriston20@gmail.com}
}

\begin{document}

\maketitle
\begin{abstract}
This paper presents the Unified Hierarchical Reasoning Controller (UHRC), a novel framework designed to enhance the reasoning capabilities of 
artificial intelligence systems. The UHRC integrates hierarchical reasoning processes, allowing for more efficient and effective decision-making in complex environments. 
We evaluate the performance of the UHRC through a series of experiments, demonstrating its superiority over traditional reasoning controllers in various tasks. 
The results indicate that the UHRC can significantly improve the adaptability and robustness of AI systems, making it a promising approach for future research in artificial intelligence and machine learning.
\end{abstract}

\section{Introduction}
\label{sec:introduction}

The autonomous control of complex nonlinear systems remains one of the most
consequential open problems in control engineering
\parencite{iqbal_nonlinear_2017, liu_new_2015, daoutidis_future_2023}.
Systems such as unmanned aerial vehicles (UAVs), autonomous ground vehicles,
and high-degree-of-freedom robotic manipulators are characterised by
high-dimensional state spaces, nonlinear inter-state couplings, actuator
saturation, and significant parametric uncertainty
\parencite{sadati_two-level_2015, ameli_robust_2021, iqbal_nonlinear_2017,
slotine_applied_1991}.
The fundamental difficulty is that the dynamical regimes governing long-horizon
navigation and those governing fast stabilisation operate on incommensurable
timescales, differ in their sensitivity to model error, and impose competing
demands on the same actuator plant
\parencite{liang_nonlinear_2018, abdelmaksoud_control_2020}.
Designing a controller that simultaneously addresses long-horizon planning and
fast reactive stabilisation across these coupled dynamics is fundamentally a
multi-scale reasoning problem, and no single architecture has proven universally
adequate across the competing demands of stability, adaptability, and
computational tractability \parencite{rojas_complex_2025, goodwin_control_2001}.
 
Classical hierarchical control has long addressed this multi-scale challenge
through explicit structural decomposition: a high-level supervisory layer
generates strategic plans and setpoints while a low-level regulatory layer
executes fast stabilisation and trajectory tracking
\parencite{singh_systems_1978, sadati_two-level_2015, scattolini_hierarchical_2007,
mesarovic_two_1970, wilson_foundations_1979}.
The theoretical underpinnings of this separation trace to the multilevel systems
framework of \parencite{mesarovic_two_1970} and the hierarchical decomposition
methods formalised in \parencite{singh_systems_1978, siljak_decentralized_2013},
which established that large-scale optimisation problems can be rendered
tractable by decomposing them into coordinated subproblems across levels.
This architecture has achieved sustained success across aerospace, process
control, and robotics, benefiting from well-established stability theory,
modular implementation, and interpretable failure modes
\parencite{sanders_hierarchical_1998, wang_hierarchical_2024, liu_generalized_2016,
bahnasawi_decentralised_1990, roberts_optimal_2001}.
The quadrotor UAV exemplifies this paradigm directly: the natural timescale
separation between slow translational dynamics and fast rotational dynamics
motivates a cascaded structure in which outer-loop position planners generate
attitude references consumed by inner-loop attitude regulators
\parencite{hoffmann_quadrotor_2007, liang_nonlinear_2018, bouabdallah_pid_2004,
nagaty_control_2013, abdelhay_modeling_2019}.
Model predictive control has provided a principled realisation of the
supervisory layer in such architectures, enabling constraint-aware planning
with formal stability guarantees at the coordination level
\parencite{scattolini_hierarchical_2007, wang_hierarchical_2024,
torrente_data-driven_2021}.
 
However, the fixed interface between layers introduces a structural information
bottleneck that becomes consequential as task complexity grows
\parencite{kolchinsky_nonlinear_2019, li_hierarchical_2023,
abdelmaksoud_control_2020}.
The high-level planner typically transmits heavily compressed information
ordinarily a kinematic waypoint sequence to the low-level regulator,
discarding contextual information regarding obstacle proximity, actuator
feasibility, and local geometric constraints that may be critical for downstream
control quality \parencite{zhu_hierarchical_2024, zhou_two-level_2017,
liu_adaptive_2013}.
Conversely, the low-level regulator executes reference trajectories without
access to the global context that motivated them, rendering local disturbance
responses globally suboptimal
\parencite{nguyen_identifying_2021, daoutidis_decomposition_2019}.
Fixed rule-based supervisory frameworks also become brittle when the system
encounters scenarios not anticipated at design time
\parencite{sadati_hierarchical_2008, dave_design_2023}.
These limitations are not incidental but structural, they arise from the
requirement that the planner and regulator communicate through a
pre-specified, fixed-dimensional interface whose information content is
determined at design time rather than adapted online to task demands
\parencite{kolchinsky_nonlinear_2019, zhang_multi-level_2015}.
 
Learning-based approaches have emerged as a principled alternative.
End-to-end reinforcement learning (RL) enables joint optimisation of planning
and regulation without manual interface specification
\parencite{almahamid_reinforcement_2021, levine_end--end_2016, kiran_deep_2022,
shakya_reinforcement_2023}.
Behavioural cloning provides a supervised alternative that sidesteps the
credit-assignment challenges of RL by learning directly from expert
demonstrations, and has seen substantial success in autonomous driving and
manipulation \parencite{levine_end--end_2016, le_mero_survey_2022, ly_learning_2021,
block_is_2024}.
However, flat end-to-end policies are notoriously sample-inefficient,
sensitive to reward shaping, and do not naturally maintain the multi-scale
temporal structure that long-horizon UAV navigation demands: the policy
allocates identical computational depth regardless of whether the task requires
strategic deliberation or reactive stabilisation
\parencite{almahamid_reinforcement_2021, agyei_deep_2025}.
Deep reinforcement learning approaches applied directly to quadrotor control
have demonstrated impressive agility in simulation
\parencite{trad_real-time_2024, saviolo_learning_2023}, but the sim-to-real gap,
sample complexity, and absence of formal stability guarantees remain open
challenges \parencite{dai_lyapunov-stable_2021, agyei_deep_2025}.
Hybrid neuro-symbolic and fuzzy-integrated controllers partially address this
brittleness by pairing learned components with structured decomposition
\parencite{tang_fuzzy_2024, mahmoud_fractional-order_2024, tran_hybrid_2021,
baheri_hierarchical_2025}, yet the supervisory-regulatory interface and its
dimensionality remain manually prescribed.
The boundary between planning and regulation is fixed by the designer, not
learned from data \parencite{borri_design_2019, yu_survey_2023}.
 
Recent advances in neural reasoning architectures offer a substantially
different approach to this problem.
Test-time scaling methods that iterate computations in latent space
\parencite{geiping_scaling_2025, li_seek_2025} and recurrent architectures that
adapt computational depth to task difficulty
\parencite{geiping_scaling_2025, li_system_2025}
suggest that the depth of processing required for a given decision can be made
emergent rather than fixed.
The Hierarchical Reasoning Model (HRM), introduced by
\parencite{wang_hierarchical_2025}, is a two-tier recurrent architecture comprising
a high-level module responsible for slow, abstract planning and a low-level
module handling rapid, detailed computation, coupled through a shared carry
state that enables adaptive allocation of computational depth.
Trained on as few as 1{,}000 examples without pre-training or chain-of-thought
supervision, the HRM achieves near-perfect accuracy on complex Sudoku problems
and optimal pathfinding in large mazes tasks on which state-of-the-art large
language model approaches fail entirely \parencite{wang_hierarchical_2025,
huang_towards_2023, feng_efficient_2025}.
Crucially, the hierarchical role specialisation between modules emerges through
training rather than being architecturally prescribed, mirroring the
representational dimensionality hierarchy observed across cortical regions in
the mammalian brain \parencite{badre_frontal_2018, huntenburg_large-scale_2018,
murray_hierarchy_2014}, where higher-order areas operate over longer intrinsic
timescales and higher-dimensional representations than primary sensory regions.
Building on this, \parencite{jin_decoupled_2025} demonstrated that decoupling
planning from execution in a hierarchical framework yields measurable gains in
multi-step reasoning tasks, further motivating the application of this design
principle to closed-loop control.
The hierarchical neuro-symbolic decision transformer of
\parencite{baheri_hierarchical_2025} similarly demonstrates that coupling a symbolic
planner with a learned low-level policy through a structured interface
outperforms both purely symbolic and purely neural baselines on sequential
decision tasks, providing additional evidence that the hierarchical
planning-execution separation confers a genuine computational advantage beyond
what either component achieves in isolation.
 
Nevertheless, all demonstrated HRM applications to date have been confined to
discrete, turn-based domains.
Whether the same unified reasoning mechanism can function as a closed-loop
controller for a continuous, real-time nonlinear dynamical systems where
control outputs must be physically continuous, observations arrive as dense
sensory streams, and temporal context must persist across thousands of control
cycles remains entirely unexplored in the literature.
Extending HRMs to such settings requires non-trivial architectural adaptations, discrete
token outputs must be replaced by continuous commands, and the supervised
training signal must be derived from a physically meaningful expert policy
rather than symbolic ground truth labels.
Localisation and mapping, which are prerequisite for any meaningful planning
computation, must be supplied by an onboard sensor
\parencite{gupta_simultaneous_2022, milijas_comparison_2021,
munguia_uav_2025, chen_slam_2022}
whose outputs form part of the continuous observation stream consumed by
the controller.
 
This paper addresses that gap by proposing and evaluating the
\textit{Unified Hierarchical Reasoning Controller} (UHRC) for the autonomous
navigation of a six-degree-of-freedom quadrotor UAV, constituting the first
systematic application of the HRM architecture to a continuous nonlinear
closed-loop control system.
The UHRC adapts the HRM's two-level reasoning structure into a recurrent
transformer controller trained via behavioural cloning
\parencite{levine_end--end_2016, le_mero_survey_2022, block_is_2024}
on expert trajectory data generated by an A* path planner
\parencite{lavalle_planning_2006, qin_review_2023}
across 500 simulated episodes.
The architecture is evaluated against a classical hierarchical baseline
incorporating SLAM-based localisation \parencite{gupta_simultaneous_2022,
milijas_comparison_2021}, A* path planning, and cascaded PID control
\parencite{bouabdallah_pid_2004, ren_cascade_2016, abdelhay_modeling_2019,
lopez-sanchez_pid_2023}
across five structured test cases of increasing environmental complexity.

The primary contributions of this work are as follows:
\begin{enumerate}
    \item The first empirical demonstration of an HRM-based architecture operating as a
          closed-loop controller for a continuous nonlinear dynamic system, establishing
          the functional feasibility of unified hierarchical reasoning in the control domain.

    \item A UHRC architecture incorporating dual-stream fusion of kinematic and LiDAR
          observations, a rolling carry buffer for persistent temporal context, Rotary
          Positional Encoding (RoPE) over the carry sequence, and a gated subgoal injection
          mechanism coupling H-module velocity planning to L-module wrench computation.

    \item A behavioural cloning training pipeline employing a versatile episode distribution,
          A*-generated expert demonstrations with injected perturbations and a
          structured loss combining weighted action imitation, velocity subgoal supervision,
          and cosine alignment.

    \item A rigorous comparative empirical evaluation quantifying the performance robustness
          trade-off between unified learned hierarchy and classical explicit decomposition
          across open-field transit, constrained corridor navigation, gap threading, dynamic
          obstacle avoidance, and parametric model uncertainty.
\end{enumerate}

\printbibliography

\end{document}
